% 2026, Prof. Dr. Matthias Jung (DL9MJ)
\begin{tikzpicture}

\begin{axis}[
    width=\linewidth*\getDarcImageFactor,
    height=0.75\linewidth*\getDarcImageFactor,
    xmin=-8.8,
    xmax=11.8,
    ymin=-3.5,
    ymax=10.8,
    axis lines=none,
    clip=false,
]

% --------------------------------------------------------------------------
% Koordinatenachsen
% --------------------------------------------------------------------------

\draw[-Triangle]
    (axis cs:0,0) -- (axis cs:11.4,0);

\draw[-Triangle]
    (axis cs:0,0) -- (axis cs:-8.5,0);

\draw[-Triangle]
    (axis cs:0,0) -- (axis cs:0,10.5);

\draw[-Triangle]
    (axis cs:0,0) -- (axis cs:0,-3.3);

% Achsenbeschriftungen

\node[anchor=west]
    at (axis cs:11.45,-0.25)
    {$U_{CE}$ in \si{\volt}};

\node[anchor=east]
    at (axis cs:-8.55,-0.25)
    {$I_B$ in \si{\micro\ampere}};

\node[anchor=south]
    at (axis cs:0.25,10.5)
    {$I_C$ in \si{\milli\ampere}};

\node[anchor=north]
    at (axis cs:0.25,-3.3)
    {$U_{BE}$ in \si{\volt}};

% --------------------------------------------------------------------------
% Skalen rechts: U_CE in Volt
% --------------------------------------------------------------------------

\pgfplotsinvokeforeach{2,4,6,8,10}{
    \draw
        (axis cs:#1,-0.10) -- (axis cs:#1,0.10);

    \node[anchor=north]
        at (axis cs:#1,-0.13)
        {\small #1};
}

% --------------------------------------------------------------------------
% Skalen links: I_B in Mikroampere
% Eine Koordinateneinheit entspricht 10 µA.
% --------------------------------------------------------------------------

\def\DARCIbTick#1/#2\relax{
    \draw
        (axis cs:#1,-0.10) -- (axis cs:#1,0.10);

    \node[anchor=north]
        at (axis cs:#1,-0.13)
        {\small #2};
}
\pgfplotsinvokeforeach{-2/20,-4/40,-6/60,-8/80}{
    \expandafter\DARCIbTick#1\relax
}

% --------------------------------------------------------------------------
% Skalen oben: I_C in Milliampere
% --------------------------------------------------------------------------

\pgfplotsinvokeforeach{2,4,6,8,10}{
    \draw
        (axis cs:-0.10,#1) -- (axis cs:0.10,#1);

    \node[anchor=east]
        at (axis cs:-0.14,#1)
        {\small #1};
}

% --------------------------------------------------------------------------
% Skalen unten: U_BE in Volt
% Die geometrische Skala ist für die gemeinsame Darstellung gestreckt.
% --------------------------------------------------------------------------

\def\DARCUbeTick#1/#2\relax{
    \draw
        (axis cs:-0.10,#1) -- (axis cs:0.10,#1);

    \node[anchor=east]
        at (axis cs:-0.14,#1)
        {\small #2};
}
\pgfplotsinvokeforeach{-1/{0{,}6},-2/{0{,}7},-3/{0{,}8}}{
    \expandafter\DARCUbeTick#1\relax
}

% --------------------------------------------------------------------------
% Stromsteuerkennlinie I_C = beta * I_B
%
% Bei beta = 100 gilt:
% 10 µA Basisstrom entsprechen ungefähr 1 mA Kollektorstrom.
% --------------------------------------------------------------------------

\addplot[
    DARCblue,
    very thick,
    domain=0:8,
    samples=2,
]
({-x},{x});

\node[
    anchor=south east,
    align=right,
]
at (axis cs:-6.4,8.2)
{Stromsteuerkennlinie\\$I_C \approx \beta I_B$};

% --------------------------------------------------------------------------
% Eingangskennlinie
%
% Vereinfachte Exponentialkennlinie der Basis-Emitter-Strecke.
% Parameter u entspricht U_BE.
% --------------------------------------------------------------------------

\addplot[
    DARCblue,
    very thick,
    domain=0.50:0.80,
    samples=160,
    variable=\u,
]
(
    {-0.054*exp((\u-0.55)/0.05)},
    {-10*(\u-0.50)}
);

\node[
    anchor=north east,
]
at (axis cs:-6.0,-3.0)
{Eingangskennlinie};

% --------------------------------------------------------------------------
% Ausgangskennlinien
%
% Vereinfachtes Modell:
% I_C nähert sich für größere U_CE dem durch I_B festgelegten Wert an.
% --------------------------------------------------------------------------

\pgfplotsinvokeforeach{1,2,3,4,5,6,7,8}{
    \addplot[
        DARCblue,
        thick,
        domain=0:10,
        samples=150,
    ]
    {#1*(1-exp(-x/0.25))};
}

\node[
    anchor=south west,
]
at (axis cs:4.5,9.5)
{Ausgangskennlinien};

% Beschriftung ausgewählter Basisströme

\def\DARCoutputLabel#1/#2\relax{
    \node[
        anchor=west,
        font=\scriptsize,
        fill=white,
        inner sep=1pt,
    ]
    at (axis cs:9.7,#1)
    {$I_B=\qty{#2}{\micro\ampere}$};
}
\pgfplotsinvokeforeach{2/20,4/40,6/60,8/80}{
    \expandafter\DARCoutputLabel#1\relax
}

% --------------------------------------------------------------------------
% Arbeitsgerade
%
% Beispiel:
% U_0 = 10 V
% R_C = 1 kOhm
%
% I_C = (U_0 - U_CE) / R_C
% Bei der gewählten Skalierung ergibt sich I_C[mA] = 10 - U_CE[V].
% --------------------------------------------------------------------------

\addplot[
    black,
    densely dashed,
    very thick,
    domain=0:10,
    samples=2,
]
{10-x};

% --------------------------------------------------------------------------
% Arbeitspunkt Q
%
% U_CE = 5 V
% I_C  = 5 mA
% I_B  = 50 µA
% --------------------------------------------------------------------------

\addplot[
    only marks,
    mark=*,
    mark size=2.3pt,
    DARCgreen,
]
coordinates {
    (5,5)
};

% --------------------------------------------------------------------------
% Hilfslinien zwischen den Kennlinien und dem Arbeitspunkt
% --------------------------------------------------------------------------

% Projektion von I_B auf die Stromsteuerkennlinie und zum Arbeitspunkt

\draw[
    DARCgreen,
    thick,
    densely dashed,
]
(axis cs:-5,0)
    -- (axis cs:-5,5)
    -- (axis cs:5,5)
    -- (axis cs:5,0);

% Zugehöriger Punkt auf der Eingangskennlinie:
% bei I_B ungefähr 50 µA ergibt sich U_BE ungefähr 0,78 V.

\draw[
    DARCgreen,
    thick,
    densely dashed,
]
(axis cs:-5,0)
    -- (axis cs:-5,-2.77)
    -- (axis cs:0,-2.77);

\fill[DARCgreen]
    (axis cs:-5,5)
    circle[radius=1.5pt];

\fill[DARCgreen]
    (axis cs:-5,-2.77)
    circle[radius=1.5pt];

% --------------------------------------------------------------------------
% Beispielhafte Aussteuerung entlang der Arbeitsgeraden
% --------------------------------------------------------------------------

\node[
    rotate=-45,
    anchor=south,
    fill=white,
    inner sep=0pt,
    outer sep=5pt,
    font=\small,
]
at (axis cs:5.0,5.0)
{\textcolor{DARCgreen}{Arbeitspunkt}};

% --------------------------------------------------------------------------
% Signalübertragung durch das Kennlinienfeld
% --------------------------------------------------------------------------

% Scheitelwerte des Eingangssignals, über die drei Kennlinien bis zum
% Ausgangssignal projiziert.
\draw[gray,densely dashed]
    (axis cs:0.75,-2.59)
    -- (axis cs:-3.5,-2.59)
    -- (axis cs:-3.5,3.5)
    -- (axis cs:6.5,3.5)
    -- (axis cs:6.5,-1.55);

\draw[gray,densely dashed]
    (axis cs:0.75,-2.90)
    -- (axis cs:-6.5,-2.90)
    -- (axis cs:-6.5,6.5)
    -- (axis cs:3.5,6.5)
    -- (axis cs:3.5,-1.55);

% Eingangssignal u_BE(t): horizontal, da U_BE senkrecht aufgetragen ist.
\addplot[
    DARCred,
    thick,
    domain=0.75:2.75,
    samples=100,
    variable=\t,
]
    ({\t},{-2.745 + 0.155*sin(720*(\t-0.75)/2.0)});

% Ausgangssignal u_CE(t): quer zur waagerechten U_CE-Achse. Die größere
% Auslenkung zeigt die Verstärkung, das Minuszeichen die Phasenumkehr.
\addplot[
    DARCred,
    very thick,
    domain=-1.55:-0.35,
    samples=100,
    variable=\t,
]
    ({5.0 - 1.5*sin(360*(\t+1.55)/1.20)},{\t});

\end{axis}

\end{tikzpicture}
